\documentclass[11pt,reqno]{amsart}
\usepackage[T1]{fontenc}
\usepackage{lmodern,amsmath,amssymb,mathtools,microtype,booktabs,array}
\usepackage[margin=1in]{geometry}
\usepackage[hidelinks]{hyperref}
\ifdefined\pdfinfoomitdate\pdfinfoomitdate=1\fi
\ifdefined\pdftrailerid\pdftrailerid{}\fi
\ifdefined\pdfsuppressptexinfo\pdfsuppressptexinfo=15\fi
\hypersetup{pdftitle={General Theta Foundations I: Acquired Geometry and Causal Resource Transfer},pdfauthor={Qian Qi}}
\numberwithin{equation}{section}
\newtheorem{theorem}{Theorem}[section]
\newtheorem{proposition}[theorem]{Proposition}
\newtheorem{lemma}[theorem]{Lemma}
\newtheorem{corollary}[theorem]{Corollary}
\theoremstyle{definition}\newtheorem{definition}[theorem]{Definition}
\newtheorem{example}[theorem]{Example}
\theoremstyle{remark}\newtheorem{remark}[theorem]{Remark}
\newcommand{\E}{\mathbb E}\newcommand{\Prb}{\mathbb P}
\newcommand{\R}{\mathbb R}\newcommand{\one}{\mathbf1}
\newcommand{\Q}{\mathcal Q}\newcommand{\TV}{\mathrm{TV}}
\newcommand{\cp}{\mathrm{cp}}\newcommand{\ev}{\mathrm{ev}}\newcommand{\st}{\mathrm{st}}
\DeclareMathOperator{\tr}{tr}\DeclareMathOperator{\dist}{dist}
\DeclareMathOperator{\law}{Law}\DeclareMathOperator{\rank}{rank}
\emergencystretch=2em
\title[Acquired geometry and causal resource transfer]{General Theta Foundations I:\\Acquired Geometry and Causal Resource Transfer}
\author{Qian Qi}\date{October 2026}
\subjclass[2020]{Primary 60G40, 62C20; Secondary 28A80, 60J05, 81P15, 94A29}
\keywords{causal experiment, predictive quotient, acquired geometry, task order, stopped prediction, finite memory, singular instrument}
\begin{document}
\begin{abstract}
We prove a two-sided resource transfer theorem for prepared causal experiments equipped with a paid reference certificate and a common acquisition program. The comparison is made on unnormalized report and task measures, along actual pre-call occupation, rather than on a postulated parameter geometry or a fictitious reference path. A single affine order combines relative calibration with accumulated absolute defects and yields implementable finite-state upper bounds and all-strategy lower bounds at the same physical budget. Relative task order alone preserves the acquired resource curve at every scale, including zero-loss strata, without the acquisition-amplified floor associated with an absolute error estimate. We characterize this order by measurable positive matrix measures on changing-rank supports, and construct Borel predictive quotients for compact-exhausted continuous-response experiments. The same theorem gives a classical controlled-channel realization and a noncommuting quantum realization with pure posterior strata; these concern post-command audits. Nonintervening directional sensing and the earlier absolute-calibration theory are retained separately in full. Event-stopped observers may reuse live labels for terminal outputs, while state-stopped observers pay for disjoint roles. Causal simulator state, stopping clocks, numerical work, calibration, raw acquisitions and physical time remain separate resources.
\end{abstract}
\maketitle
\section{Introduction}
An experiment does not acquire a manifold. It acquires reports through preparations and interventions, each of which has a causal position and a physical cost. The future tests executable after those reports determine which distinctions matter. Only after this experiment and its tests have been fixed does a geometric question arise: how much of the resulting predictive information can be retained and used with finite resources?

The central difficulty is not the definition of a predictive quotient. It is the passage from a quotient and an acquired law to an algorithm that uses the permitted interface, and to a converse that applies to all competing algorithms. Three obstacles occur together. Selection changes the acquired law and must pay for rejected observations. An error harmless for an ordinary report can be amplified by a rare stopping event. Finally, a label used to route a live protocol need not be unavailable as a terminal output. A geometric argument that omits any of these features can give an incorrect resource curve.

We address these obstacles by comparing \emph{unnormalized task measures}. Given a genuinely prepared joint report--audit law $K_0(dx,dz)$, let $\nu$ be its report law and derive
\[
 f(x)=\E_0[Z\mid x],\qquad w(x)=\E_0[\|Z-f(x)\|^2\mid x].
\]
Neither quantity is prescribed as an information metric. For a dictionary $C$ the task distortion is $w+\dist(f,C)^2$. A probability measure selected at expected paid-attempt budget $b$ is subject to a density cap, leading to the profile
\begin{equation}\label{eq:intro-profile}
 \Q_m(b)=\min_{|C|\le m,\ 0\le h\le b,\ \int h\,d\nu=1}
       \int [w+\dist(f,C)^2]h\,d\nu.
\end{equation}
This is a marked population trimmed-quantization functional. Trimming and finite dictionaries have a classical statistical history \cite{CGM97,GG99}; the claim here is not the invention of that functional. The problem is to prove when its geometry describes a paid, causal, unknown-model experiment.

Theorem~\ref{thm:transfer} is the central result. One actual operating model supplies a converse against \emph{every} legal strategy. A separate, parameter-independent acquisition program supplies an upper bound for \emph{every} operating model and every surviving history. We compare only the current productive report--audit row while retaining the actual pre-call occupation of the competitor. This yields an accepted comparison measure, not a supposed execution under the comparison experiment. Its mass, density cap and task loss can consequently be controlled without transferring a stopping time to an experiment in which the rule may never terminate.

The new feature is the simultaneous treatment of order and defect. Relative order of nonnegative task measures changes constants and effective budgets. Absolute defects are integrated against actual pre-call occupation. In particular, when there is relative order but no absolute defect, fixed nonzero calibration distortion does not impose an additive risk floor. Corollary~\ref{cor:uniform} gives matching resource curves at all scales. This differs from an estimate whose error grows as acquisition count times an absolute one-call discrepancy. Both regimes are consequences of the same theorem, rather than unrelated approximation results.

The order has a raw verification principle. The positive matrix measure
\[
 J(D)=\int_{D\times H}(1,z)(1,z)^\mathsf{T}\,P(dy,dz)
\]
contains exactly the localized affine-square tests. Proposition~\ref{prop:matrix-order} identifies the order of these tests with the order of $J$. Theorem~\ref{thm:singular-order} expresses it on the measurable range of $J$, including variable rank and zero-mass reports. It does not invert a vanishing eigenvalue uniformly. The construction is task-derived; it neither assumes a Fisher matrix nor identifies the matrix rank with a memory exponent.

There are two further structural conclusions. Theorem~\ref{thm:quotient} constructs the Borel predictive quotient and a measurable section from compact exhaustion and continuous response coordinates. It replaces an assumed Borel quotient realization by verifiable sufficient conditions; it does not assert that every measurable equivalence relation has such a realization. Lemma~\ref{lem:roles} gives separate exact packing rules for event and state stopping. The distinction persists under the resource-certified simulator composition of Proposition~\ref{prop:composition}.

The geometric consequences depend on the actual small-ball mass of the prepared law and on global covering. They permit singular location measures, a variable-rank task carrier and a positive acquired zero-loss atom. Two raw realizations verify the same general theorem: finite classical controlled channels and finite-dimensional noncommuting quantum instruments. Their targets are post-command audits. Neither is an estimation theorem for an untouched pre-intervention state. The full preceding theory, including nonintervening directional Gaussian sensing, absolute-calibration occupation bounds, finite-description examples and stopped-flow realization, is supplied as the complete technical companion \cite{Companion}. It is preserved, not relabelled as new mathematics.

\subsection*{Relation to experiment comparison and predictive representations}
Classical experiment comparison \cite{Blackwell} and predictive-state minimality \cite{SC} provide important antecedents. They do not by themselves give the particular paid stopping profile, a finite-state strategy lift, or a separate simulator clock and workspace account. Our theorem uses a reference/effort witness and a common task-order construction. It is not a necessary-and-sufficient classification of all causal strategies. The exact order characterization below is for a declared family of localized affine-square tests, not for the existence of arbitrary history-dependent balancing programs. These distinctions are part of the mathematical statement, not qualifications added after the proof.

\section{Prepared kernels, acquired task laws and observers}\label{sec:data}
All hidden-state, action and report spaces are standard Borel. A fixed unknown index $\theta\in\Theta$ specifies an actual preparation and positive normalized causal kernels. Positivity means nonnegative measures, not an everywhere strictly positive density. Dependence on the entire physical past may be included in the hidden state. A legal policy sees past visible reports and its admitted randomization, but not $\theta$ or a future audit. Iteration of these kernels defines its finite-prefix laws; the consistent countable path law is their kernel extension.

Each call has a cost $c_i\ge c_*>0$. The decision to make the next call and its productive flag are predictable. Commitment is permitted only after a productive report, and occurs exactly once almost surely in each model. At that time a specified future audit $Z\in H$ is scored, where $H\subset\R^k$ is nonempty, compact and convex. The prediction is committed before that audit is observed. A prospective joint report--audit kernel is defined at every productive call, including rejected calls. The stopping coin is independent of the audit conditional on the entire current visible report and the past. Write $L=\operatorname{diam}(H)^2$; the case $L=0$ is trivial whenever a legal rule exists.

The core cost $C_\tau$ includes initial preparation charged to the core, failed calls and all productive calls through commitment. A separately executed final audit must be added to the raw budget. This convention is fixed before taking a risk infimum. Physical time, stored program constants and temporary computation are not identified with call count.

\subsection*{Three resource interfaces}
A checkpoint procedure may use its whole precommit history, but must output one element of a fixed dictionary of at most $M$ values. This counts its output width only; it is not a claim of $\log_2 M$ total working bits.

An autonomous observer has an alphabet $S$, $|S|\le M$, and one fixed parameter-independent program. At an intercall boundary its only persistent observer memory is its label $s\in S$. Its next command distribution is a fixed function of $s$. A physical port name, elapsed time, subprogram index, random seed or former report is not an extra free input. During a call its current command, current report and fresh coins are available. Whatever persists to another call must be encoded in $S$. A fixed decoder $r:S\to H$ reads only the committed label.

In the \emph{event-stopped} interface a post-report kernel emits $(s',e)\in S\times\{0,1\}$. An event $e=1$ ends the experiment and retains $s'$ for the decoder; no subsequent command is executed. This public commitment event is not a stored bit used in continuing computation. The same label may be used live on one path and terminal on another. In the \emph{state-stopped} interface a fixed map $h_*:S\to\{0,1\}$ determines halt status. The initial label is live, and type-one labels are globally absorbing terminal labels with no executed command rows. Entry into a terminal label is allowed only after a productive report. Thus live and terminal types are disjoint in this interface, and only in this interface by definition.

Write $R_i^\Theta(N,M)$, $i\in\{\cp,\ev,\st\}$, for the infimum of worst-case expected squared loss over programs with robust almost-sure commitment and $\sup_\theta\E_\theta C_\tau\le N$. An empty legal class has value $+\infty$. A program chosen once at random is not a free persistent seed. The ideal measurable class may have fixed real constants; effective implementation is stated separately in Section~\ref{sec:composition}.

\subsection*{A physically acquired comparison law}
Execute a declared prepared comparison experiment and obtain its joint report--audit probability $K_0(dx,dz)$ on $X\times H$. Its report marginal is $\nu$. The law is genuinely preparable and its specification has its own resource account. It need not be a member of $\Theta$. In particular no finite set of calibration data is silently promoted to a uniform certificate. Set
\begin{equation}\label{eq:derived}
 f(x)=\int zK_0(dz\mid x),\quad
 w(x)=\int\|z-f(x)\|^2K_0(dz\mid x),\quad
 d_r(x)=w(x)+\|f(x)-r\|^2.
\end{equation}
Disintegration gives Borel versions with $f\in H$; null sets may be completed arbitrarily. The profile is
\begin{equation}\label{eq:profile}
 \Q_m(b)=\min_{\substack{C\subset H,\ 1\le|C|\le m\\
                  0\le h\le b,\ \int h\,d\nu=1}}
        \int [w+\dist(f,C)^2]h\,d\nu,
 \qquad b\ge1,\ m\ge1.
\end{equation}
We use $+\infty$ for an infeasible cap or width. The quantities $f,w$ are task coordinates of acquired predictive information, not a complete substitute for the predictive quotient. Their derivation and the quotient construction in Appendix~\ref{sec:quotient} precede any use of local dimension.

\subsection*{Lower certificates on actual histories}\label{sec:lower-cert}
Fix \emph{one} actual model $\theta_*\in\Theta$, independent of the competing program. At each actual reached productive pre-call history $h$ and legal command $a$, write $P_{h,a}(dy,dz)$ for its conditional full current report--audit law. Let $P^0_{h,a}$ be a measurable comparison row on the same report space, with a Borel mark $x=\pi_{h,a}(y)$ of law $\nu$ and
\begin{equation}\label{eq:full-moments}
 \E^0[Z\mid y]=f(\pi_{h,a}y),\qquad
 \E^0[\|Z\|^2\mid y]=w(\pi_{h,a}y)+\|f(\pi_{h,a}y)\|^2.
\end{equation}
These are conditional identities given the \emph{full} report, not just its mark. Versions and inequalities are required on actual reachable histories, not only almost everywhere under an unrelated comparison path law.

Fix $0<l_-\le1\le u_-<\infty$. For all Borel $q:Y\to[0,1]$ and all finite-image Borel $r:Y\to H$, suppose predictable nonnegative defects $a_-(h,a),e_-(h,a)$ obey
\begin{align}
 \int q\,d(P_{h,a})_Y
 &\le u_-\int q\,d(P^0_{h,a})_Y+a_-(h,a),\label{eq:lower-mass}\\
 \int q(y)\|z-r(y)\|^2P_{h,a}(dy,dz)
 &\ge l_-\int q(y)\|z-r(y)\|^2P^0_{h,a}(dy,dz)-e_-(h,a).
 \label{eq:lower-task}
\end{align}
Independent current randomization can be adjoined to both rows: integrating inequalities for each fixed random value proves the same inequalities with the same defects. All command-conditioned lower inequalities hold before averaging the chosen command.

Let $A_\tau$ count productive calls. At this same actual model require, for every legal competing procedure,
\begin{equation}\label{eq:effort}
 \E_* C_\tau\ge\kappa_-\E_* A_\tau,\qquad\kappa_->0.
\end{equation}
At budget $N$ suppose, uniformly over all robustly feasible competitors,
\begin{equation}\label{eq:occupation}
 \E_*\sum_{j\le\tau:\ j\ {
m productive}}a_-(H_{j-1},a_j)\le A_N,
 \qquad
 \E_*\sum_{j\le\tau:\ j\ {
m productive}}e_-(H_{j-1},a_j)\le E_N.
\end{equation}
These are actual occupation certificates, not properties merely of the proposed upper program. Uniform row defects give $A_N\le(N/\kappa_-)\bar a_-$ and $E_N\le(N/\kappa_-)\bar e_-$. Lemma~\ref{lem:drift} supplies a nonuniform alternative.

\subsection*{One common upper construction}\label{sec:upper-cert}
Independently, let $\mathcal G$ be a parameter-independent program which, from each of its start histories, reaches one productive report. Its conditional expected attempt cost is at most $\kappa_+>0$ on every surviving history in every $\theta\in\Theta$. It uses $B$ live roles and returns to its start label on rejection; no reset of the physical state is assumed. It forms a virtual mark $x$, a current choice $J$ and a permitted export $T_Jr$ of its task readout. For every finite dictionary $C$ assume
\begin{equation}\label{eq:exports}
 |\{T_jr:j\in J,\ r\in C\}|\le D_0|C|,
 \qquad D_0\in\mathbb N.
\end{equation}
Every continuing phase and coin needed for this construction is included in the $B$ live roles. For fixed $0<l_+\le1\le u_+<\infty$ and $a_+,e_+\ge0$, assume for all Borel $q:X\to[0,1]$ and finite-image Borel $r:X\to H$, conditionally on every such history,
\begin{align}
 \E[q(x)]&\ge l_+\int q\,d\nu-a_+,\label{eq:upper-mass}\\
 \E[q(x)\|Z-T_Jr(x)\|^2]
 &\le u_+\int q(x)d_{r(x)}(x)\nu(dx)+e_+.
 \label{eq:upper-task}
\end{align}
The expectation here includes the within-attempt choice $J$. Unlike the lower, the upper is not required to hold for each choice separately. The choices retained after this call are nevertheless charged through the actual export or live role. These hypotheses express one common causal construction, not modelwise optimizers.

\section{The causal task-order transfer theorem}\label{sec:central}
Suppose that every terminating state-stopped competitor reaches at least $B_0\ge1$ live labels under $\theta_*$. Such a certificate must concern actual physical command compatibility, not printed action names. Define
\begin{equation}\label{eq:widths}
\begin{array}{c|cc}
 i&m_-(M)&m_+(M)\\\hline
 \cp&M&\lfloor M/D_0\rfloor\\
 \ev&M&\lfloor M/D_0\rfloor\\
 \st&M-B_0&\lfloor(M-B)/D_0\rfloor .
\end{array}
\end{equation}
The event upper also requires $M\ge B$. The trivial state-stopped lower $B_0=1$ is valid for a procedure initially waiting to make a call.

\begin{theorem}[Affine task-order transfer with a paid reference and common construction]\label{thm:transfer}
Assume all the prepared-kernel, full-report, task-order, effort, occupation, resource and common-construction hypotheses of Section~\ref{sec:data}. Fix $N,M$, put $b_-=N/\kappa_-$, $b_+=N/\kappa_+$, and suppose $b_-\ge1$ and $m_-(M)\ge1$. Set
\[
 t_N=(1-A_N)_+/u_-.
\]
If $t_N>0$, every one of the three risks satisfies
\begin{equation}\label{eq:central-lower}
 R_i^\Theta(N,M)\ge
 \bigl[l_-t_N\Q_{m_-(M)}(b_-/t_N)-E_N\bigr]_+.
\end{equation}
If $t_N=0$, the asserted lower is zero, with no evaluation of $\Q$ at infinity. An empty legal class still has risk $+\infty$.

Suppose also $m_+(M)\ge1$, $b_+\ge1$, and, for event stopping, $M\ge B$. If
\begin{equation}\label{eq:beta}
 \beta_N=\frac{l_+b_+}{1+b_+a_+}\ge1,
\end{equation}
one common Borel program with actual expected cost at most $N$ in every model attains
\begin{equation}\label{eq:central-upper}
 R_i^\Theta(N,M)\le
 \frac{u_+(1+b_+a_+)}{l_+}\Q_{m_+(M)}(\beta_N)+b_+e_+.
\end{equation}
It commits almost surely. If $\beta_N<1$, the always-accept first-attempt construction instead gives $R_i^\Theta(N,M)\le u_+\Q_{m_+(M)}(1)+e_+$ at cost at most $\kappa_+\le N$.

The lower is against every legal competing strategy and uses one actual witness. The upper is one parameter-independent strategy for all models. The two effort constants need not agree. Neither direction assumes that the comparison experiment belongs to $\Theta$ or that a normalized predictive update is uniformly contracting.
\end{theorem}

\begin{corollary}[Scale-uniform relative transfer]\label{cor:relative}
Under Theorem~\ref{thm:transfer}, suppose $A_N=E_N=a_+=e_+=0$. If $l_+b_+\ge1$, then
\begin{equation}\label{eq:relative}
 \frac{l_-}{u_-}\Q_{m_-(M)}(u_-b_-)
 \le R_i^\Theta(N,M)
 \le\frac{u_+}{l_+}\Q_{m_+(M)}(l_+b_+).
\end{equation}
The factors $l_\pm,u_\pm$ may be fixed nonunit constants, independent of budget and resolution. In particular no term proportional to $b_\pm$ times a fixed absolute calibration error is required in this relative-order class.
\end{corollary}
\begin{proof}
Substitute $t_N=1/u_-$ and $\beta_N=l_+b_+$ in the theorem.
\end{proof}

\begin{corollary}[Separate exact checkpoint and autonomous laws]\label{cor:exact}
Suppose the certificates are exact: $l_\pm=u_\pm=1$, all defects vanish, $\kappa_-=\kappa_+=\kappa$, and exports are the identity. Then, for $N\ge\kappa$,
\[
 R_\cp^\Theta(N,m)=\Q_m(N/\kappa).
\]
The constructed autonomous upper uses $B+m$ state-stopped labels, or $\max(B,m)$ event-stopped labels. If the mandatory live count is sharp, $B_0=B$, then
\begin{align}
 R_\st^\Theta(N,B+m)&=\Q_m(N/\kappa)\quad(m\ge1),\label{eq:typed-exact}\\
 R_\ev^\Theta(N,M)&=\Q_M(N/\kappa)\quad(M\ge B).\label{eq:event-exact}
\end{align}
The subtraction by $B$ belongs only to the state-stopped model. Nonidentity exports require the expanded dictionary in \eqref{eq:widths}, unless an optimal dictionary is invariant under all exports. These exact stopping distinctions are inherited from the complete companion and are reproved here to fix the resource semantics.
\end{corollary}
\begin{proof}
The same lower and upper profile coincide in each asserted interface. Identity exports give $D_0=1$. The packing and converse widths are proved in Lemma~\ref{lem:roles}. No lower on live states is used to forbid event-stopped reuse.
\end{proof}

The familiar absolute-defect theorem is also a specialization: take $l_\pm=u_\pm=1$ while retaining $A_N,E_N,a_+,e_+$. Thus relative and absolute calibration are composed in one statement. The chosen metric is the common scored task throughout; it is not changed between the converse and the construction.

\section{Proof of the central theorem}
\begin{lemma}[Profile attainment]\label{lem:attainment}
For every finite $m\ge1$ and $b\ge1$, the minimum in \eqref{eq:profile} is attained. The profile is nonincreasing in $m,b$ and left-continuous in $b>1$.
\end{lemma}
\begin{proof}
Write a dictionary as an ordered tuple in the compact set $H^m$, permitting repetitions. The feasible densities are a weak-star closed subset of the radius-$b$ ball of $L^\infty(\nu)$: positivity, the upper bound and the mass equality are each closed under pairing with $L^1(\nu)$. They are therefore weak-star compact. As the tuple varies, its squared-distance function varies uniformly on $H$, since all points belong to a bounded set. Along a convergent product net, split the objective difference into this uniform variation and the weak-star pairing against the fixed bounded measurable loss at the limit tuple. This proves continuity of the objective and attainment. Nets avoid an unwarranted sequential compactness assumption. Monotonicity follows from inclusion of feasible sets. For left continuity take an optimizer $h$ at $b>1$ and let $h_\lambda=(1-\lambda)h+\lambda$. It has unit mass and cap $b_\lambda=(1-\lambda)b+\lambda<b$. Its objective tends to the optimum as $\lambda\downarrow0$, while monotonicity bounds the optimum at $b_\lambda$ between the optimum at $b$ and this objective.
\end{proof}

\begin{lemma}[Actual-history comparison measure]\label{lem:comparison}
Fix any robustly admissible competitor of cost at most $N$, with its fixed terminal dictionary $C$. There are finite nonnegative measures $\xi_c$ on $X$, $c\in C$, such that $\xi=\sum_c\xi_c$ obeys
\begin{align}
 \xi&\le (\E_*A_\tau)\nu\le b_-\nu,\label{eq:measure-cap}\\
 1&\le u_-\xi(X)+A_N,\label{eq:measure-mass}\\
 \E_*\|Z-\widehat r\|^2
 &\ge l_-\sum_{c\in C}\int d_c\,d\xi_c-E_N.
 \label{eq:measure-risk}
\end{align}
The measure $\xi$ is not asserted to be the terminal law of any comparison execution.
\end{lemma}
\begin{proof}
Adjoin an independent uniform random variable to the current report on both sides, sufficient to represent the competitor's randomized output and stopping decision. A finite-output kernel has such a representation by cumulative probabilities. At a fixed reached productive history $h$ and chosen command $a$, the post-report decision then gives an acceptance indicator $q_h(y)$ and an output $r_h(y)\in C$ on acceptance. Give $r_h$ an arbitrary fixed value on rejection. The notation $y$ now includes the current independent random variable. Integrating the original certificates over that variable preserves their constants. Dependence on the entire earlier history is permitted in these functions.

For each Borel $D\subset X$, set
\begin{equation}\label{eq:xi-definition}
 \xi_c(D)=\E_*\sum_{j\le\tau:\ j\ {
m productive}}
 \int\one_{\{\pi_{H_{j-1},a_j}(y)\in D\}}
 q_{H_{j-1}}(y)\one_{\{r_{H_{j-1}}(y)=c\}}
 P^0_{H_{j-1},a_j}(dy,H).
\end{equation}
The outer expectation is the actual pre-call occupation at $\theta_*$. All terms are jointly measurable and nonnegative. Tonelli defines the measures, whose total is finite because $\E_*A_\tau\le b_-$. Summing over $c$ and dropping the acceptance indicator gives \eqref{eq:measure-cap}: each comparison row has mark law $\nu$.

The analogous sum with the actual rows instead of $P^0$ has mass exactly one, because the actual procedure commits exactly once almost surely. Apply \eqref{eq:lower-mass} at each actual reached history, then sum. The occupation certificate yields \eqref{eq:measure-mass}. Apply \eqref{eq:lower-task} in the same way to the disjoint actual commitment events. Their expected losses sum to the actual risk. The comparison contribution is $d_{r_h(y)}(\pi_{h,a}y)$ by \eqref{eq:full-moments}; this identity conditions on the full current report and therefore does not overlook extra information seen by the competitor. The accumulated loss defect is at most $E_N$, giving \eqref{eq:measure-risk}. One may first truncate the number of calls. Monotone convergence for the nonnegative contributions and the finite defect bounds justify the limit.

No actual past is replaced by a comparison past. In particular the competitor need not terminate, or obey a cost bound, if hypothetically run under the comparison experiment. This distinction is why the argument applies when that experiment is not an operating model.
\end{proof}

\begin{lemma}[Live roles, output roles and physical compatibility]\label{lem:roles}
A program with $\ell$ live roles and $t$ output roles has a state-stopped implementation with $\ell+t$ labels and an event-stopped implementation with $\max(\ell,t)$ labels, preserving every path cost and complete stopped scored law. Conversely an event-stopped $M$-label program has at most $M$ decoder values. A state-stopped $M$-label program with at least $B_0$ reached live labels has at most $M-B_0$ terminal decoder values.

For a physical live-state certificate, consider finitely many mandatory pre-call contexts $I$ with no free context input. Suppose the commands legal at a history of context $i$ are exactly $\mathcal U_i$, and every terminating competitor visits every context with positive probability under the same witness. Let $b_{\rm phys}$ be the least size of a cover of $I$ by subsets $J$ with $\bigcap_{i\in J}\mathcal U_i\ne\varnothing$. Then at least $b_{\rm phys}$ live labels are required. This certificate is invariant under duplicate names and bijective stateless recodings of the same physical commands.
\end{lemma}
\begin{proof}
For state stopping take disjoint copies of the live and terminal role sets and type them zero and one. For event stopping embed the live and terminal role sets separately into one alphabet of size $\max(\ell,t)$. On continuing transitions use the live embedding; on commitment emit the final event and use the terminal embedding. A command row is read only on a continuing branch, while the decoder is read only after commitment. Sharing their label therefore causes no conflict and uses no continuing mode bit. Couple the original and recoded coins. Every finite stopped prefix and every path cost agree, and almost-sure termination gives equality of the stopped scored laws. A fixed decoder has one value per retained label. In state stopping the global halt map excludes each reached live label from the output set, proving its subtraction. Event stopping has no such exclusion.

A live label reached in contexts $J$ must have one command distribution supported on all the corresponding $\mathcal U_i$. Since $J$ is finite, the union of the exceptional null command sets is null; hence their intersection supports a probability and is nonempty. The sets of contexts served by the labels cover $I$, proving the bound. Conversely a nonempty intersection contains a command and hence a Dirac command distribution. Replacing names without changing executed kernels, legal domains, costs and reports does not change these intersections. A translator which itself remembers a hidden phase is not such a recoding; its memory must be charged.
\end{proof}

\begin{proof}[Proof of Theorem~\ref{thm:transfer}]
Fix a competing program. Lemma~\ref{lem:roles} bounds its decoder width by $m_-(M)$. Let $s=\xi(X)$ in Lemma~\ref{lem:comparison}. If $t_N>0$, then \eqref{eq:measure-mass} gives $s\ge t_N$. Scale the measure to $\xi'=(t_N/s)\xi$. It has mass $t_N$ and remains bounded above by $b_-\nu$. Thus $h=(d\xi'/d\nu)/t_N$ has unit mass and cap $b_-/t_N$. Nearest-dictionary projection and nonnegativity imply
\[
 \sum_{c\in C}\int d_c\,d\xi_c
 \ge\int [w+\dist(f,C)^2]\,d\xi
 \ge t_N\Q_{m_-(M)}(b_-/t_N).
\]
Together with \eqref{eq:measure-risk} this proves the asserted lower for this competitor. Its worst-model risk dominates its risk at $\theta_*$. Taking the infimum over competitors gives \eqref{eq:central-lower}, without exchanging an infimum and a supremum. Nonnegative loss gives the zero truncation and the stated case $t_N=0$.

For the upper assume \eqref{eq:beta}. Choose an attained optimizer $(C,h)$ at width $m_+(M)$ and cap $\beta_N$. Let $r_C(x)$ be the first nearest point to $f(x)$ in a fixed ordering of $C$. After each attempt of $\mathcal G$, accept $x$ with probability $h(x)/\beta_N$, export $T_Jr_C(x)$ and commit; otherwise return to the start routing label. All these operations are Borel. The exports are assumed jointly measurable and take at most $D_0|C|$ distinct values, so their finite decoder and transitions are Borel as well.

Conditional on every surviving attempt history, the acceptance probability is at least
\begin{equation}\label{eq:hazard}
 p=\frac{l_+}{\beta_N}-a_+=\frac1{b_+}>0.
\end{equation}
Consequently the probability of reaching attempt $j+1$ is at most $(1-p)^j$, by induction on conditional survival probabilities. In particular commitment is almost sure and the expected number of attempts is at most $1/p=b_+$. Each begun attempt has conditional expected cost at most $\kappa_+$, not merely a cost conditional on a favorable outcome. Summing the predictable attempt costs gives actual expected cost at most $\kappa_+b_+=N$.

The conditional unnormalized loss contributed by an attempt is at most $u_+\Q_{m_+(M)}(\beta_N)/\beta_N+e_+$. Sum its contributions on the disjoint terminal events, using the same survival bound. The result is at most
\[
 b_+\left[\frac{u_+}{\beta_N}\Q_{m_+(M)}(\beta_N)+e_+\right],
\]
which is \eqref{eq:central-upper}. The program has $B$ live roles and at most $D_0m_+(M)$ terminal roles. Lemma~\ref{lem:roles} gives exactly the sufficient widths in \eqref{eq:widths}. A checkpoint counts only the terminal width. No physical renewal law was used: the conditional hazard controls survival even if the hidden state after rejection differs arbitrarily from its former state.

If $\beta_N<1$, commit after the first attempt, and use an optimizer at cap one. The loss certificate with $q=1$ gives $u_+\Q_{m_+(M)}(1)+e_+$, and the whole-attempt cost is at most $\kappa_+$. This proves all remaining assertions.
\end{proof}

\section{Raw task order and changing-rank stability}\label{sec:order}
The central theorem is formulated in measures because selection acts before normalization. This section gives a verification principle from primitive kernels and a singular extension of that principle. It does not assume a covariance or information matrix in the experiment.

\subsection*{A measure determined by executable affine-square tests}
For a report--audit law $P(dy,dz)$ on $Y\times H$ define the finite positive matrix measure
\begin{equation}\label{eq:moment-matrix}
 J_P(D)=\int_{D\times H}v(z)v(z)^\mathsf T P(dy,dz),
 \qquad v(z)=(1,z_1,\ldots,z_k)^\mathsf T.
\end{equation}
For finite positive matrix measures write $A\preccurlyeq B$ when $B(D)-A(D)$ is positive semidefinite for every Borel $D$. This definition is independent of the choice of dominating scalar measure.

\begin{proposition}[An exact order criterion for localized affine-square tests]\label{prop:matrix-order}
For two joint laws $P,P^0$ on the same full report space and $0<l\le1\le u$, the following assertions are equivalent:
\begin{equation}\label{eq:matrix-order}
 lJ_{P^0}\preccurlyeq J_P\preccurlyeq uJ_{P^0};
\end{equation}
for every Borel report event $D$ and every $a\in\R^{k+1}$,
\begin{equation}\label{eq:affine-tests}
 l\int_D|a^\mathsf T v(z)|^2dP^0
 \le\int_D|a^\mathsf T v(z)|^2dP
 \le u\int_D|a^\mathsf T v(z)|^2dP^0.
\end{equation}
It suffices to test rational vectors $a$. These equivalent conditions imply, for all Borel $q\in[0,1]$ and finite-image Borel $r:Y\to H$, both the report-mass and the squared-task inequalities with zero additive defect and these same factors. A violation of \eqref{eq:matrix-order} is witnessed by a single Borel report event and a rational affine-square test.
\end{proposition}
\begin{proof}
The integral in \eqref{eq:affine-tests} is $a^\mathsf T J_P(D)a$. The definition of positive semidefiniteness gives the equivalence. Continuity of the quadratic form gives sufficiency of rational vectors and a rational witness for a strict violation. Taking $a=(1,0,\ldots,0)$ gives mass order. For a fixed readout $r$, sum the $k$ quadratic tests with vectors $(-r_j,e_j)$, obtaining the order of $\|z-r\|^2$. Partition $Y$ into the finitely many cells where $r$ is constant. For a simple nonnegative $q$, apply order on intersections with these cells and add. Bounded monotone approximation extends to arbitrary $q\in[0,1]$. All measures are finite and all tested losses bounded.
\end{proof}

Thus the criterion is necessary and sufficient for the declared task-rich affine-square family. It is a sufficient condition for the narrower squared-prediction task of the central theorem. It is not claimed necessary for that narrower task, or for the existence of an arbitrary causal balancing strategy. The matrix is built from raw joint audit laws, not from an assumed parameterization. The upper construction may use its virtual report space after averaging its current command; the lower criterion must use the full report seen by a competitor, conditional on that competitor's chosen command.

\begin{theorem}[Order on measurable singular ranges]\label{thm:singular-order}
Let $J_0,J_1$ be finite positive matrix measures of size $k+1$. With
$\lambda=\tr J_0+\tr J_1$, take measurable positive semidefinite densities $A_0,A_1$ with respect to $\lambda$. Then
\begin{equation}\label{eq:singular-equivalence}
 lJ_0\preccurlyeq J_1\preccurlyeq uJ_0
\end{equation}
is equivalent, $\lambda$-almost everywhere, to equality of kernels
$\ker A_0=\ker A_1$ and to
\begin{equation}\label{eq:range-spectrum}
 lP_0\preccurlyeq A_0^{\dagger/2}A_1A_0^{\dagger/2}
            \preccurlyeq uP_0,
\end{equation}
where $P_0$ is the orthogonal projection onto $\operatorname{ran}A_0$ and $\dagger$ denotes the Moore--Penrose inverse. On the zero range both sides are zero. All these maps are measurable, even when rank changes. No positive lower bound on a nonzero eigenvalue of $A_0$ is required. In particular every zero affine-square task on an acquired report event remains exactly zero under the order.
\end{theorem}
\begin{proof}
A finite signed matrix measure is positive if and only if its density is positive semidefinite almost everywhere. To see the nontrivial direction, test the density against all rational vectors; a negative set of positive measure for one such quadratic form contradicts positivity of its integral on that set. The countable intersection of the resulting full-measure sets handles all rational vectors, hence all real vectors by continuity. Thus \eqref{eq:singular-equivalence} is equivalent to $lA_0\preccurlyeq A_1\preccurlyeq uA_0$ almost everywhere.

For two such matrices, a vector in $\ker A_0$ has zero quadratic form under $A_1$, and positivity implies that it is in $\ker A_1$. The lower inequality gives the converse. On their common range $A_0$ is positive definite; congruence by its inverse square root yields \eqref{eq:range-spectrum}. Conversely equality of kernels makes $A_1$ supported on that range. Congruence of \eqref{eq:range-spectrum} by $A_0^{1/2}$ gives the original two matrix inequalities, extended by zero on the common kernel.

Measurability follows from finite-dimensional spectral functional calculus: $P_0$ is the pointwise limit of $A_0(A_0+n^{-1}I)^{-1}$, and the scalar function equal to $x^{-1/2}$ on $x>0$ and zero at zero is Borel. Applying it spectrally gives a Borel matrix function on each rank stratum, hence on their finite measurable union. We use these matrices pointwise; their unbounded inverse norms are never taken as a uniform error constant. Finally a nonnegative task of zero integral under $J_0$ has integral at most $u$ times zero under $J_1$, proving preservation of zero tasks.
\end{proof}

There is an important distinction between this assertion and a rank formula for risk. Rank is allowed to change; the rate still comes from the acquired masses and coverings of Section~\ref{sec:geometry}. For example a nonzero matrix density with arbitrarily small positive eigenvalues can obey the order with fixed $l,u$, while an arbitrarily small absolute perturbation can create a new positive direction in its kernel and fail the order. This is precisely the difference relevant to a zero-loss stratum.

\begin{proposition}[Positive-kernel verification and relative command calibration]\label{prop:positive-order}
Suppose raw positive input states satisfy $ls_0\le s\le us_0$ in an ordered vector space and a fixed positive normalized report--audit instrument sends an input to a joint probability law. Its output laws and task matrices satisfy the same order. More generally let $\mathcal I_j$ be positive instruments and let a nominal probability vector $(\lambda_j^0)$ satisfy a desired joint-law identity after averaging. If
\[
 (1-\eta)\lambda_j^0\le\lambda_j\le(1+\eta)\lambda_j^0
 \quad\hbox{for every }j,\qquad 0\le\eta<1,
\]
then the actual averaged joint law lies between $1-\eta$ and $1+\eta$ times the nominal law, for every input. Successive relative input and instrument bounds multiply their lower factors and upper factors. They do not add an acquisition-amplified defect.
\end{proposition}
\begin{proof}
Positivity sends the two positive differences $s-ls_0$ and $us_0-s$ to positive measures. Integration of $v(z)v(z)^\mathsf T$ against a positive measure preserves positive matrix order. The command assertion follows by multiplying each nonnegative instrument measure by its nonnegative weight difference and summing. Applying the two inequalities successively multiplies the factors. These arguments do not divide by any report probability.
\end{proof}

\subsection*{Actual-mass-weighted predictive stability}
\begin{proposition}[Continuity through likelihood zeros and pure states]\label{prop:continuity}
A classical stochastic instrument on signed measures, or a finite-dimensional completely positive normalized quantum instrument, maps normalized inputs $s,t$ to unnormalized post-report states $A_y,B_y$. Use any common finite dominating measure for these two outputs and let $p_y=u(A_y)$, $q_y=u(B_y)$ be their masses. Then
\begin{align}
 \int\|A_y-B_y\|_1dy&\le\|s-t\|_1,\label{eq:stable1}\\
 \|p-q\|_{\TV}&\le\tfrac12\|s-t\|_1,\label{eq:stable2}\\
 \int\min(p_y,q_y)\left\|\frac{A_y}{p_y}-\frac{B_y}{q_y}\right\|_1dy
 &\le2\|s-t\|_1.\label{eq:stable3}
\end{align}
Terms of zero mass are zero. The same unnormalized and report-law inequalities hold for a complete stopped instrument under one common legal policy which terminates almost surely under both inputs.
\end{proposition}
\begin{proof}
Write $s-t=D_+-D_-$ with positive parts. Positivity, the triangle inequality and preservation of integrated mass give
$\int\|\mathcal I_y(D_+)-\mathcal I_y(D_-)\|_1dy\le u(D_+)+u(D_-)=\|s-t\|_1$.
Taking masses gives \eqref{eq:stable2}. If $p\le q$, then
\[
 p\|A/p-B/q\|_1\le\|A-B\|_1+(q-p).
\]
Interchanging the roles for $q\le p$, integrating, and using $|p-q|\le\|A-B\|_1$ proves \eqref{eq:stable3}. A finite causal continuation with a fixed policy is again a positive normalized instrument. Truncate a stopped continuation at a deterministic horizon, retaining a distinct surviving outcome. Apply the finite inequality, then let the horizon increase. The surviving masses vanish under both inputs by almost-sure termination, and the sum over terminal times converges monotonically. No expected-time assertion is deduced from these norm inequalities.
\end{proof}

The upper finite-state construction never stores a normalized posterior. It uses the acquired report and a finite decoder. Consequently stability of the raw instrument, rather than a globally contracting posterior recursion, suffices for the singular realizations below. More general closed-loop optimal-control claims would require their own strategy comparison and stability; they are not consequences of comparing the same fixed policy.

\section{Acquired geometry at regular, stratified and atomic scales}\label{sec:geometry}
All masses in this section refer to the actually prepared pushforward
$\sigma=(f,\sqrt w)_\#\nu$. They are not volumes on a parameter space. Conditional expectation of any feasible density given $(f,\sqrt w)$ preserves its mass, density cap and objective for every dictionary; conversely a density on this pushforward lifts by composition. Thus the profile can equivalently be optimized on this actual marked law. The transfer theorem compares operating task measures with this physical acquisition law. In the relative-order realizations it also bounds actual upper-program report measures above and below by fixed multiples of $\nu$. Thus the comparison remains tied to acquired mass, even though perturbed conditional means need not coincide with $f$.

\begin{proposition}[A non-power geometric comparison]\label{prop:sublevel}
For $r\ge0$ let
\[
 \Omega_m(r)=\max_{C\subset H,\ 1\le|C|\le m}
       \nu\{x:w(x)+\dist(f(x),C)^2\le r^2\}.
\]
For all $b,m\ge1$,
\begin{equation}\label{eq:sublevel-bounds}
 \sup_{r\ge0}r^2[1-b\Omega_m(r)]_+
 \le\Q_m(b)
 \le\inf\{r^2:\Omega_m(r)\ge1/b\}.
\end{equation}
In particular a verified estimate $b\Omega_m(r_-)\le1/2$ and a verified covering by $m$ distortion sublevel sets of total acquired mass at least $1/b$ at radius $r_+$ imply $r_-^2/2\le\Q_m(b)\le r_+^2$. The assertion requires no integer local dimension or isotropic geometry.
\end{proposition}
\begin{proof}
For a fixed dictionary and a density cap $b$, at most $b\Omega_m(r)$ of the selected mass lies in its radius-$r$ sublevel set. The remaining mass costs at least $r^2$, proving the lower. To see that the maximum defining $\Omega_m$ is attained, use an ordered tuple in $H^m$. For every report, the indicator of the closed sublevel condition is upper semicontinuous in that tuple. Reverse Fatou for these bounded indicators makes its integral upper semicontinuous. Compactness gives a maximizing tuple. If the corresponding sublevel set has mass $v\ge1/b$, the density $\one_D/v$ is feasible and has loss at most $r^2$. Taking the infimum over such radii gives the upper. This proof keeps the density cap and actual mass; a cover of the formal reachable set alone provides neither.
\end{proof}

The proposition can be used with anisotropic boxes, different local scales, or non-power small-ball laws, provided both the upper mass estimate and the required acquired cover are verified. It is a finite-scale statement, not a claim that an unspecified anisotropic dimension determines risk. We give explicit closed curves for a useful class and its finite stratifications.

Assume $\sigma$ is supported on $K\times[0,1]$, where $K\subset H$ is compact, and for some $d,\alpha>0$, $0<r_0\le1$ and $c,C>0$,
\begin{equation}\label{eq:rectangle}
 c r^d s^\alpha\le\sigma(B(x,r)\times[0,s])\le C r^d s^\alpha
 \quad(x\in K,\ 0<r\le r_0,\ 0<s\le1).
\end{equation}
These are global estimates on the stated component, including small reliability mass. Singular Ahlfors-regular location measures of noninteger $d$ are permitted.

\begin{proposition}[The marked acquired profile]\label{prop:regular}
Under \eqref{eq:rectangle}, uniformly for $b\ge1$ and integer $m\ge1$,
\begin{equation}\label{eq:F}
 \Q_m(b)\asymp F(b,m):=
 \max\{b^{-2/\alpha},(bm)^{-2/(d+\alpha)}\}.
\end{equation}
The constants depend only on the displayed data and the bounded task carrier. This geometric calculation is inherited from the stopped-experiment and predictable-geometry work retained in \cite{Companion}; its proof is included for independence of the central theorem.
\end{proposition}
\begin{proof}
Cover $K$ by finitely many balls of radius $r_0/2$. Their upper bounds yield $\sigma(K\times[0,s])\le C_0s^\alpha$. Under density cap $b$, choose $s=c_0b^{-1/\alpha}$ sufficiently small so that at most half the selected mass has reliability at most $s$. This gives a risk lower $s^2/2$.

For a dictionary with $m$ centers, its distortion sublevel at radius $r$ lies in the union of $m$ location balls of radius $r$ with reliability at most $r$. A center outside $K$ whose ball intersects $K$ can, for this mass bound, be replaced by one intersection point and the radius doubled. For $2r\le r_0$, the union has acquired mass at most $C_1m r^{d+\alpha}$. Set $r=c_1(bm)^{-1/(d+\alpha)}$ with one sufficiently small constant $c_1$ so that the restriction and the half-mass bound both hold for all $b,m\ge1$. Proposition~\ref{prop:sublevel} gives the other lower in \eqref{eq:F}.

For the upper, the location marginal has two-sided $d$-dimensional ball bounds by taking $s=1$. A maximal $4r$-separated set covers $K$ by $4r$-balls. The marginal upper bound then forces at least $c_2r^{-d}$ points for small $r$. Its radius-$r$ balls are disjoint. Choose constants $c_3>0$ small and $C_2$ large such that, for all sufficiently large $b$, the integer
\[
 j=\min\{m,\lfloor c_3b^{d/\alpha}\rfloor\}\ge1,
 \qquad r=C_2(bj)^{-1/(d+\alpha)}
\]
satisfies $r\le\min(r_0/8,1)$, $j\le c_2r^{-d}$, and $cj r^{d+\alpha}\ge1/b$. The last condition fixes a sufficiently large $C_2$; after fixing it, a sufficiently small $c_3$ enforces the packing condition because
$j r^d=C_2^d j^{\alpha/(d+\alpha)}b^{-d/(d+\alpha)}$.
Increasing the lower threshold for $b$ then enforces the radius condition and $j\ge1$. Use $j$ separated centers. The disjoint rectangles $B(x_i,r)\times[0,r]$ have union mass at least $1/b$ and distortion to those centers at most $2r^2$. Selecting a unit mass under cap $b$ from that union proves the upper, and substitution of $j$ gives \eqref{eq:F}. For $b$ in the remaining bounded interval, always accept and use one center. Its bounded loss is controlled by a fixed multiple of $b^{-2/\alpha}$, uniformly over arbitrarily large $m$. This proves all claimed uniformities.
\end{proof}

\begin{corollary}[All-scale resource curves under fixed relative calibration]\label{cor:uniform}
Suppose \eqref{eq:rectangle} and the relative hypotheses of Corollary~\ref{cor:relative} hold. On a set of feasible budgets suppose the positive effort ratio $\kappa_+/\kappa_-$ and positive width ratio $m_-/m_+$ are bounded above and below by fixed constants. Then
\begin{equation}\label{eq:uniform-law}
 R_i^\Theta(N,M)\asymp
 F(N/\kappa_-,m_-(M))
\end{equation}
throughout that set, including arbitrarily fine resolution. The constants may depend on $l_\pm,u_\pm$, but these factors need not approach one as $N,M$ grow.

More generally the same conclusion holds with absolute defects if $A_N\le A_0<1$, $b_+a_+\le A_1<\infty$, $E_N\le c_0F(b_-,m_-)$ and $b_+e_+\le C_0F(b_-,m_-)$, for a sufficiently small $c_0>0$ depending on the fixed geometric and order constants.
\end{corollary}
\begin{proof}
For pure relative order and $l_+b_+\ge1$, both effective budgets and widths in \eqref{eq:relative} differ from $b_-,m_-$ by bounded positive factors. Apply Proposition~\ref{prop:regular}. If $1\le b_+<1/l_+$, then $b_-$ lies in a bounded interval. The first-attempt upper is bounded, while the lower and $F(b_-,m_-)$ are bounded below by a fixed positive constant independently of $m_-$. This treats the boundary regime. With absolute defects, $t_N$ is bounded below by $(1-A_0)/u_-$ and the corrected upper budget differs by fixed factors. The geometric lower therefore exceeds $cF$ before subtracting $E_N$. Take $c_0\le c/2$. The upper defect has the stated order, yielding the result. A fixed positive absolute defect need not satisfy these scale conditions; relative order and absolute proximity are distinct hypotheses.
\end{proof}

\begin{theorem}[A mixed-dimensional acquired law]\label{thm:stratified}
Suppose the acquired marked law decomposes as a finite probability mixture
$\sigma=\sum_{j=1}^J\lambda_j\sigma_j$, with $\lambda_j>0$, $\sum_j\lambda_j=1$. The components are distributions of the actual task coordinates, not the pushforwards of unspecified hidden conditional means. Each $\sigma_j$ satisfies \eqref{eq:rectangle} with its own $d_j,\alpha_j$ and fixed geometric constants. Define the feasible acquired-mass simplex
\[
 \mathcal S_b=\{(s_j):s_j\ge0,\ \sum_js_j=1,\ s_j\le b\lambda_j\}.
\]
For $m\ge2J$ and all $b\ge1$,
\begin{equation}\label{eq:stratified-law}
 \Q_m(b)\asymp\min_{s\in\mathcal S_b}
 \sum_{j:s_j>0}s_j
 \max\left\{
 \left(\frac{b\lambda_j}{s_j}\right)^{-2/\alpha_j},
 \left(\frac{b\lambda_jm}{s_j}\right)^{-2/(d_j+\alpha_j)}\right\}.
\end{equation}
A zero $s_j$ contributes zero. The constants depend on the finite component data and $J$, not on $b,m$. Applying Theorem~\ref{thm:transfer} to this profile gives the corresponding two-sided causal resource law; no stratum is renormalized without retaining its selected mass $s_j$.
\end{theorem}
\begin{proof}
For any feasible density and dictionary, set $s_j=\lambda_j\int h\,d\sigma_j$. The masses belong to $\mathcal S_b$. When $s_j>0$, the normalized density $\lambda_jh/s_j$ on the $j$th component is bounded by $b\lambda_j/s_j$. Its dictionary still has at most $m$ elements. Thus the full objective is at least
\[
 \sum_{j:s_j>0}s_j\Q_m^j(b\lambda_j/s_j).
\]
This proves the lower after Proposition~\ref{prop:regular}.

Conversely fix feasible $s$ and allocate $m_j=\lfloor m/J\rfloor$ centers to every component with positive mass. Their total is at most $m$ and $m_j\ge m/(2J)$. Choose an optimal component density $h_j$ at cap $b\lambda_j/s_j$ and dictionary $C_j$. The measure $\zeta=\sum_{s_j>0}s_jh_j\sigma_j$ has mass one and is bounded by $b\sigma$. Use its Radon--Nikodym density with respect to $\sigma$ and the union of the $C_j$. Nearest-center distortion to that union is no greater than distortion to the component's own dictionary. This gives the upper $\sum_js_j\Q_{m_j}^j(b\lambda_j/s_j)$, within a fixed factor of the expression in \eqref{eq:stratified-law}. If a component label is not observed, this construction still uses only the density of $\zeta$ on the actual marked coordinates; no latent component choice is made by the observer. The displayed objective is continuous at $s_j=0$ by its positive powers of $s_j$, and the feasible simplex is compact and nonempty. Its minimum is attained. Taking this minimum completes the proof.
\end{proof}

The mixture theorem permits different local dimensions and reliability exponents, including noninteger singular dimensions. It is more than a statement that strata may be supplied through an interface: it gives their coupled acquisition--memory optimization and proves both sides with the original masses. Proposition~\ref{prop:sublevel} remains available when an anisotropic component has no single power-law description.

\begin{proposition}[An acquired zero-loss atom]\label{prop:atom}
Suppose $\sigma=\lambda\delta_{(x_*,0)}+(1-\lambda)\sigma_0$, $0<\lambda<1$, where $\sigma_0$ has no mass at that atom. Then $\Q_m(b)=0$ for every $m\ge1$ if $b\lambda\ge1$. If $s=1-b\lambda>0$, then
\begin{equation}\label{eq:atom}
 s\Q_m^0\left(\frac{b(1-\lambda)}s\right)
 \le\Q_m(b)\le
 s\Q_{m-1}^0\left(\frac{b(1-\lambda)}s\right)\quad(m\ge2).
\end{equation}
For $m=1$ only the displayed lower applies; a valid upper fixes its center at $x_*$ and optimizes the remaining selected mass with that center. If $w>0$ almost surely on the remainder, then every finite-width profile is positive below threshold. Under pure relative transfer, zero risk is impossible when $u_-b_-\lambda<1$, and is attainable when $l_+b_+\lambda\ge1$ and the upper width is feasible. In the exact matched case the threshold is exactly $N=\kappa/\lambda$.
\end{proposition}
\begin{proof}
Above threshold accept only the atom and decode $x_*$. Below threshold at least mass $s$ must come from the remainder. Scale any selected remainder measure down to mass $s$; it remains dominated by $b(1-\lambda)\sigma_0$ and nonnegative loss cannot increase. Normalize it to obtain the lower in \eqref{eq:atom}. For the upper retain all available atom mass $b\lambda$, allocate one center to it, and optimize the mass-$s$ remainder with $m-1$ centers. With one center there is no zero-center optimization, which explains the separate statement. If the remainder has $w>0$ almost surely, any positive selected probability has positive residual loss; attainment from Lemma~\ref{lem:attainment} makes the profile minimum strictly positive. Finally substitute the atom criterion into the two effective caps of \eqref{eq:relative}. Relative calibration preserves zero tasks but may change the required acquisition threshold by the stated factors; it does not preserve the numerical threshold without exact matching.
\end{proof}

\subsection*{Nonuniform defects and metastable preparation}
\begin{lemma}[Killed-potential occupation bound]\label{lem:drift}
Let a nonnegative predictable defect $d_i$ be charged at raw call $i$, and suppose, for every legal command before commitment, nonnegative integrable potentials obey
\[
 d_i+\E_*[V_i\one_{\{i<\tau\}}\mid\mathcal F_{i-1},a_i]
 \le V_{i-1}+\gamma\E_*[c_i\mid\mathcal F_{i-1},a_i].
\]
Here $V_0$ is integrable, $\gamma\ge0$, and after commitment the potential is killed. Then every budget-$N$ competitor satisfies $\E_*\sum_{i\le\tau}d_i\le\E_*V_0+\gamma N$. Potentials may depend on the hidden physical state; this is an analytic certificate, not uncharged observer memory.
\end{lemma}
\begin{proof}
Multiply by the predictable indicator of making call $i$, integrate and sum to a deterministic horizon $n$. All interior potentials cancel. The remaining final potential is nonnegative and may be dropped. Bound the expected partial cost by $N$, then use monotone convergence for the nonnegative defect sum. Randomized commands are handled by integrating the commandwise inequality. No optional stopping with unverified uniform integrability is needed.
\end{proof}

\begin{corollary}[Relative order with rare defective modes]\label{cor:metastable}
Suppose relative order holds on good raw histories, while visible conditional additive defects are bounded by the conditional expectations of $a_b\one_{\{D_{i-1}=1\}}$ and $e_b\one_{\{D_{i-1}=1\}}$ for a hidden mode $D\in\{0,1\}$. At every legal call, given the physical past, assume
\[
 \Prb_*(D_i=1\mid D_{i-1}=0,\cdot)\le u,
 \qquad\Prb_*(D_i=0\mid D_{i-1}=1,\cdot)\ge v>0.
\]
Let $\pi_0=\Prb_*(D_0=1)$ be the actual prepared mass and $c_i\ge c_*>0$. The central theorem applies with
\[
 A_N=\frac{a_b}{v}\pi_0+\frac{a_bu}{vc_*}N,
 \qquad E_N=\frac{e_b}{v}\pi_0+\frac{e_bu}{vc_*}N.
\]
These quantities may replace rowwise worst-case defects even if bad-mode errors are not small and the residence time $1/v$ is large.
\end{corollary}
\begin{proof}
Use $V(D)=(a_b/v)\one_{\{D=1\}}$. In a good mode its expected increase is at most $a_bu/v$; in a bad mode its expected decrease is at least $a_b$ and pays for the defect. Killing at commitment only improves the inequality. Since each raw call costs at least $c_*$, Lemma~\ref{lem:drift} gives the first bound. Counting only productive calls makes it smaller. Replace $a_b$ by $e_b$ for the task bound and use conditional expectation to pass from physical modes to the visible-history defects. This retains actual preparation and selection-biased residence, rather than substituting an invariant-mode distribution.
\end{proof}

\section{Classical and quantum realizations from raw kernels}\label{sec:realizations}
The two realizations below are verified after the general theorem. Recurrent witness preparation, when present, belongs to the charged productive primitive; it is not a reset made freely available to other models. Neither supplies a premise used to prove it. Their conditional physical evolution between calls may be noncontracting, and their acquired laws may be singular. Both tasks score the post-command audit. A target which must remain unchanged is a different experiment; the nonintervening directional result is retained in the companion.

\subsection*{A common paid physical interface}
The gate port $G$ costs one per call and has success probability at least $p_0>0$, conditional on every physical and visible past. A success enables exactly one productive call at port $A$, also of cost one. Commitment is forbidden at $G$ and allowed after a productive report. Rejection at $A$ returns the protocol to $G$; it need not reset the hidden state. The two ports have disjoint physical command domains and no port tag is a free intercall input to an autonomous observer. Suppose the actual lower witness has gate success probability at most $p_1\le1$. Then
\begin{equation}\label{eq:gate-constants}
 \kappa_-=1+1/p_1,\qquad \kappa_+=1+1/p_0,
 \qquad B_0=B=2.
\end{equation}
Indeed every productive call consumes a gate success. A terminating path has no unconsumed success. Predictable counting gives $\E_*A_\tau\le p_1\E_*G_\tau$, so $\E_*C_\tau\ge\kappa_-\E_*A_\tau$. The probability of $j$ failures after any gate start is at most $(1-p_0)^j$, giving conditional expected gate count at most $1/p_0$ and whole-attempt cost at most $\kappa_+$. Both mandatory ports are visited before commitment, and Lemma~\ref{lem:roles} gives two live roles. Consequently the construction requires at least two event labels or three state labels when it has one output. If the witness has constant success $p_0=p_1$, then $\E_*C_\tau=\kappa\E_*A_\tau$ with $\kappa=1+1/p_0$. In that matched case budget feasibility is exactly $N\ge\kappa$ with the respective state thresholds. Even when $p_0=1$, the cost is two, not one.

\begin{theorem}[A relatively calibrated classical realization]\label{thm:classical}
Let a hidden symbol $E$ take values in a finite set of size $q$. A finite group $\mathcal G$ acts transitively on that set. The allowed productive commands apply $g\in\mathcal G$ and then, within the same call, a fixed positive normalized joint report--audit kernel $J_{gE}(dy,dz)$ with bounded audit in $H$. There is no unknown intervening operation between that command and its audit instrument. Evolution between port interactions is otherwise arbitrary and causal.

The genuinely prepared comparison has uniform input and joint law $K_0=q^{-1}\sum_eJ_e$. Its full report is its mark. Suppose the operating family contains a single witness whose conditional input weights, after every reached pre-command history, satisfy
\[
 (1-\eta_r)/q\le\Prb_*(E=e\mid h)\le(1+\eta_r)/q,
 \qquad 0\le\eta_r<1.
\]
Suppose one parameter-independent implementation chooses its command independently of the hidden current symbol conditional on its visible past, with conditional weights
\[
 (1-\eta_u)/|\mathcal G|\le\lambda_g\le(1+\eta_u)/|\mathcal G|,
 \qquad0\le\eta_u<1,
\]
uniformly over all models and surviving histories. Use the gate interface above. Then Theorem~\ref{thm:transfer} holds with identity exports, zero additive defects, and
\begin{equation}\label{eq:classical-relative}
 l_-=1-\eta_r,\quad u_-=1+\eta_r,
 \qquad l_+=1-\eta_u,\quad u_+=1+\eta_u.
\end{equation}
The comparison need not be an operating member. Fixed nonzero $\eta_r,\eta_u<1$ preserve the all-scale resource curve of Corollary~\ref{cor:uniform} whenever its acquired geometry holds.
\end{theorem}
\begin{proof}
A chosen permutation preserves the two componentwise bounds on the witness weights. Multiplying by each positive $J_e$ and summing gives the full command-conditioned joint-law order
$(1-\eta_r)K_0\le P_{h,g}\le(1+\eta_r)K_0$.
Take $P^0_{h,g}=K_0$ and the identity mark. Conditional moments given its full report are exactly \eqref{eq:derived}, so all lower certificates follow from Proposition~\ref{prop:matrix-order}. This holds for every chosen command and every competing policy, not merely for a randomizing policy.

For any fixed incoming symbol, uniform averaging over a transitive group gives the uniform distribution of its image: each image has the same number of preimages in the group. Averaging also over any incoming distribution preserves that fact. The actual componentwise command bounds and positivity therefore put the full joint upper-program law between $1-\eta_u$ and $1+\eta_u$ times $K_0$. Independence of the command from the hidden current input is needed for this conditional averaging; it is stated as a physical implementation property, not inferred from marginal coin frequencies. The upper accepts on the resulting report, not on a discarded group choice. Exports are the identity because the audit is the post-command audit. Group selection is within-call work, and only the gate or productive routing label or final decoder value persists. The gate proof establishes the same-witness cost lower, common cost upper and both live roles. The general theorem now applies.
\end{proof}

For a concrete acquired geometry take $q=2$, with identity and sign reversal. Independently draw $(Y,T)$ from a fixed law on $K\times[0,1]$, report these coordinates, and audit
\begin{equation}\label{eq:classical-mark}
 Z=Y+TEe_1
\end{equation}
in a sufficiently large compact convex carrier. The uniform preparation gives $f=Y,w=T^2$. Choose a $d$-regular location measure and $\Prb(T\le s)=s^\alpha$ to verify \eqref{eq:rectangle} directly; choose a finite mixture of such laws to verify Theorem~\ref{thm:stratified}; add a physically prepared point mass at $(x_*,0)$ to obtain Proposition~\ref{prop:atom}. A biased fresh witness and imperfect command probabilities satisfy \eqref{eq:classical-relative} with fixed nonzero constants. On $T=0$ the task is exactly deterministic even under these errors. On $T>0$ the two-point audit moment matrix has rank two, while at $T=0$ it has rank one. This is an actual changing-rank task measure, not a formal limiting parameter set.

The command group and its bit-level sampler are part of the program description and workspace account. A high-cardinality group does not provide free randomization or a free persistent seed. In the binary realization a single fresh sign coin suffices within a productive call. Conditional posterior continuity of the classical raw kernels is governed by Proposition~\ref{prop:continuity}, including likelihood zeros.

\begin{theorem}[A relatively calibrated noncommuting quantum realization]\label{thm:quantum}
Let the productive physical system have finite dimension $d\ge2$. Its available commands include the $d^2$ Weyl conjugations, followed within the same productive call by a fixed positive normalized quantum report--audit instrument $\mathcal J$, with classical joint law $\mathcal J(\rho)(dy,dz)$ and $z\in H$. The physical task is again the post-command audit. Put $\omega=I/d$ and $K_0=\mathcal J(\omega)$.

Suppose one actual witness obeys, conditionally before every selected productive command,
\[
 (1-\eta_r)\omega\preccurlyeq\rho\preccurlyeq(1+\eta_r)\omega,
 \qquad0\le\eta_r<1.
\]
Suppose the common implementation uses current Weyl weights between $(1-\eta_u)/d^2$ and $(1+\eta_u)/d^2$, independently of the current hidden input conditional on the past, with $0\le\eta_u<1$. With the same gate certificates, all conclusions of Theorem~\ref{thm:classical} hold with the constants \eqref{eq:classical-relative}. The operating witness need not be exactly maximally mixed. Incoming states for the upper may be pure, and post-report ranks may change on positive acquired mass. No common eigenbasis or normalized-filter contraction is assumed.
\end{theorem}
\begin{proof}
Let $W_{p,q}$ be the shift by $q$ followed by the phase $\exp(2\pi ipj/d)$ on the $j$th basis vector. On a matrix unit $|j\rangle\langle k|$, averaging the phase kills off-diagonal terms; averaging shifts makes the diagonal uniform. Hence
\[
 d^{-2}\sum_{p,q}W_{p,q}\rho W_{p,q}^*=I/d
\]
for every density matrix. This elementary randomization identity is standard \cite{AMTW}; we use its complete positive-instrument consequence, not a new state-randomization claim.

Each chosen conjugation fixes $\omega$ and preserves the witness Loewner order. Positivity of the full report--audit instrument therefore gives lower joint-law order relative to $K_0$ with factors $1\pm\eta_r$. Positivity means that the image of every positive operator is a positive classical measure, including both report and future audit. It consequently compares the full observable report and task, not just the report marginal.

For the upper, the componentwise weight bounds give
\[
 (1-\eta_u)\omega
 \preccurlyeq\sum_{p,q}\lambda_{p,q}W_{p,q}\rho W_{p,q}^*
 \preccurlyeq(1+\eta_u)\omega
\]
for every incoming state, including pure states. Apply the same positive instrument. Proposition~\ref{prop:matrix-order} supplies the task certificates, and identity export, counted live roles and the gate costs finish the verification. The lower witness has a relative input bound around $\omega$; we do not describe that hypothesis as absence of an input-rank condition. The absence of a floor concerns the post-report states and task strata, where the following example has genuine rank changes.
\end{proof}

Here is a raw qubit instrument giving the same acquired geometry as \eqref{eq:classical-mark} but noncommuting predictive states. Draw $(y,t)$ with any stated law $\nu_0$ on $K\times[0,1]$, choose an equatorial unit vector $n=(\cos\phi,\sin\phi,0)$ and $0\le\gamma(y,t)\le1$, and set
\[
 E_v=(I+v\gamma n\cdot\sigma)/2,\qquad L_v=\sqrt{E_v},\qquad v=\pm1.
\]
The report is $(y,t,v)$ and the unnormalized post-report state is $L_v\rho L_v$. At commitment measure $\sigma_z$, obtain $e=\pm1$, and score $Z=y+te e_1$. Under $\omega=I/2$, each $v$ has mass $1/2$ and the normalized posterior is $E_v$. Its $z$ component is zero, so the future audit sign is fair. Thus the reference report law is $\nu_0$ times a fair $v$, and $f=y,w=t^2$, derived directly from the instrument.

The eigenvalues of $E_v$ are $(1\pm\gamma)/2$. Taking $\gamma=1$ on a positive-mass set and $\gamma<1$ on another gives actual pure and full-rank posterior strata. Varying $\phi$ produces noncommuting effects. For a continuous compact realization one may take $\gamma(t)=1$ for $t\le1/2$ and $\gamma(t)=3/2-t$ for $t\ge1/2$, and a continuous nonconstant angle. The two-sided task order remains uniform across these strata; Proposition~\ref{prop:continuity} supplies actual-mass-weighted predictive-update and report continuity. A point mass at $t=0$ also gives a changing-rank task matrix and an exact zero-loss acquired component.

The qubit is part of the physical experiment, not a quantum memory register owned by the finite classical observer. The pulse and instrument form one productive call in this convention; their duration, finite sampling precision and arithmetic are separately charged. An interface requiring a separate pulse call must add its cost and any intervening routing state. Both realizations retain a strictly positive paid gate floor $p_0$.

\subsection*{Predictive realization rather than an assumed parameter manifold}
For the finite classical kernels and finite-dimensional quantum instruments whose entries depend continuously on compact report coordinates, the primitive continuity criterion in Proposition~\ref{prop:raw-quotient} supplies compact positive-evidence history patches. Their executable predictive quotient is therefore Borel by Theorem~\ref{thm:quotient}. Fixed continuous causal evolutions between calls are included and need not contract. For a completely arbitrary merely Borel hidden evolution, the risk theorem still uses the stated raw certificates, but we do not assert the compact-exhaustion quotient theorem without its regularity hypotheses. The unknown model's family of posteriors is never stored by the common upper program.

\section{Finite computation, calibration and causal composition}\label{sec:composition}
The inequalities proved so far concern Borel existence with a declared physical acquisition program. A finite source description and an executable precision certificate are additional mathematical data. We state exactly where they enter.

\begin{proposition}[Numerical implementation at a restored physical budget]\label{prop:numerics}
Suppose a chosen ideal density and dictionary at cap $\beta$ have a finite implementation whose acceptance probability, as a function of the current virtual report, differs pointwise from $h/\beta$ by at most $\eta\ge0$. Suppose its actual exported decoder values differ from the ideal exported values by at most $\delta_{\rm num}$, with no larger output width. These finite operations use the same legal acquisition program and counted live roles. Put $b=N/\kappa_+$ and
\[
 \beta=\frac{l_+b}{1+b(a_++\eta)}\ge1.
\]
If the chosen ideal pair has objective at most $\Q_m(\beta)+\xi$, the implemented program terminates almost surely at actual expected core cost at most $N$ and risk at most
\begin{equation}\label{eq:numerical-law}
 \left[
 \sqrt{\frac{u_+(1+b(a_++\eta))}{l_+}
                  (\Q_m(\beta)+\xi)+b(e_++L\eta)}
       +\delta_{\rm num}\right]^2.
\end{equation}
Here $\xi$ is a certified design/integration suboptimality, not a local optimizer's unverified gap. A uniform bound $T_{\mathrm{attempt}}$ on conditional expected physical and arithmetic time of a whole attempt yields expected time at most $bT_{\mathrm{attempt}}$. Program bits and scratch are separately those of the certified implementation.
\end{proposition}
\begin{proof}
Acceptance rounding reduces the conditional hazard by at most $\eta$. The chosen cap gives $l_+/\beta-a_+-\eta=1/b$. The same conditional-survival summation as in the main theorem yields at most $b$ attempts and cost at most $N$. Before perturbing decoder values, the one-attempt unnormalized loss changes by at most $L\eta$ under the actual joint law, because loss is bounded by $L$. Its other contribution is at most $u_+(\Q_m(\beta)+\xi)/\beta+e_+$. Summation gives the expression under the square root. At the actual committed report, the implemented and ideal exported values are within $\delta_{\rm num}$; Minkowski's inequality in the actual stopped $L^2$ law adds this distance to root-mean-square loss. It does not multiply the decoder error by the waiting time. The time claim is a separate predictable sum of the stated whole-attempt bound.
\end{proof}

In the preceding display $T_{\mathrm{attempt}}$ denotes a declared time quantity, not a free clock input to the observer. An implementation based on $m$ fixed centers in a $k$-dimensional bounded box, $p$-bit coordinate certificates and sequential squared-distance comparisons has sufficient read-only data $O(mkp)$ bits, within-call scratch $O(kp+\log m)$ bits and $O(mkp^2)$ bit operations per report, in addition to its primitive measurement and coordinate-evaluation costs. These are upper accounts, not optimal computational lower bounds. General Borel classifiers need not have such certificates. The complete dyadic implementation in the companion gives explicit rational thresholds, separate read-only/scratch counts and an independently proved profile lower.


\begin{corollary}[Finite-description matching upper under effective primitives]\label{cor:effective}
In the pure relative regular class, suppose additionally that the actual upper-program report laws are bounded above by $\bar u\nu$ on every surviving history. Assume effective compactness of the dictionary carrier and certified finite approximations of exported readouts inside $H$. Assume the calibrated report law has an effective sequence of finite physical partitions which approximates bounded finite decision vectors in $L^1(\nu)$, with certified integration of their masses and bounded task losses. Primitive classifiers and their finite costs are supplied as part of this effective interface. Then a finite-description upper with the same live and terminal role counts has risk at most a fixed multiple of $F(b,m)$ at actual expected cost at most $\kappa_+b$, for all feasible $b,m$. Every such construction has finite, separately charged program, scratch and execution accounts. No polynomial-time discovery or matching bit-complexity lower is asserted.
\end{corollary}
\begin{proof}
For $b\ge4/l_+$ choose the ideal cap $\beta=l_+b/2$. Its conditional acceptance is at least $2/b$. By Proposition~\ref{prop:regular}, its profile is at most a fixed multiple of $F(b,m)$. Represent its finite decision as functions $q_c(x)$, one for each dictionary value, whose sum is the acceptance probability. Approximate the vector on one finite physical partition by rational vectors in the subprobability simplex. The assumed effective approximation makes the sum of its $L^1(\nu)$ coordinate errors arbitrarily small, say at most $\gamma$. Domination of the actual report law bounds its actual conditional decision discrepancy by $\bar u\gamma$ on every history, and its unnormalized loss discrepancy by $L\bar u\gamma$. Approximate the actual finite exported values inside $H$ as well.

Choose $\gamma$ so small that $\bar u\gamma<1/(2b)$ and $bL\bar u\gamma\le F(b,m)$, and choose export error at most $\sqrt{F(b,m)}$. The conditional acceptance still exceeds $1/b$, so the expected cost is at most $\kappa_+b$. Summing loss and applying the stopped $L^2$ readout estimate gives a fixed multiple of $F(b,m)$. For $1\le b<4/l_+$, accept the first attempt and use one finitely described value in $H$. Bounded loss is a fixed multiple of the strictly positive lower bound for $b^{-2/\alpha}$ on this interval, uniformly in $m$.

The existence argument can be effected by finite enumeration using the supplied integration certificates and effective dictionary nets: the above choices leave strict mass and loss margins, so a candidate with certified acceptable intervals is eventually found. This need not be an efficient search. Current partition indices and rational randomization are charged within-call work; they do not create new intercall roles. All chosen descriptions and precision parameters are finite. Exact zero risk on an atom requires an exactly implementable atom classifier and readout; otherwise Proposition~\ref{prop:numerics}, not an unqualified finite-code zero claim, governs that component.
\end{proof}

\begin{proposition}[Resource-certified causal morphisms]\label{prop:composition}
A causal simulator must specify its parameter-independent strategy lift, legal command map, report transformation, internal state, clock and terminal task readout. Suppose a target program has $\ell$ live and $t$ terminal roles. A compatible simulator has at most $K$ continuing internal states, including every persistent phase, coin and clock, and at most $E$ terminal export classes sufficient for its final decoder. Then a source implementation uses at most $\ell K+tE$ state-stopped labels or $\max(\ell K,tE)$ event-stopped labels. Without a terminal sufficient-class certificate, retaining the whole simulator state gives the safe product bound.

If the complete stopped scored transcript, including mapped stopping time and premature commitment, is within total variation $\varepsilon$ of the target law uniformly over the declared models and policies, then the bounded common-task risk changes by at most $L\varepsilon$. Two compatible simulators multiply their continuing and terminal resource factors and add their deficiencies. Raw-call count and physical-time transfer require separate actual cost or killed-potential certificates. Reverse risk inequalities require a reverse simulator or an actual lower witness; terminal-law proximity alone does not provide them.
\end{proposition}
\begin{proof}
Before commitment retain the pair of target live role and continuing simulator state. After commitment retain the pair of target output role and sufficient export class. Compatibility supplies legal transitions and commands without an uncharged port or phase. Lemma~\ref{lem:roles} gives the two packing counts. Serial composition retains triples; rebracketing yields the products $K_1K_2$ and $E_1E_2$. For a loss in $[0,L]$, integration of its superlevel probabilities gives $|\E_P\ell-\E_Q\ell|\le L\|P-Q\|_{\TV}$. Total variation contracts under a common Markov kernel, so the triangle inequality adds complete-law deficiencies. Costs are generally unbounded functions of stopped paths and are not controlled by this bounded-loss argument; sum their certified increments or use Lemma~\ref{lem:drift} instead.
\end{proof}

Combining Propositions~\ref{prop:numerics} and \ref{prop:composition} adds $L\varepsilon$ to \eqref{eq:numerical-law} while charging the simulator's states and actual cost map. This composes approximation and simulation in one task metric. A bound on an allowed error is not, without attained alternatives, an unavoidable lower bound. The following construction supplies such alternatives on a single task.

\begin{theorem}[A same-task resource--calibration--deficiency law]\label{thm:joint}
Take any core class satisfying Theorem~\ref{thm:transfer} with identity exports and $B$ common live roles. Add one compulsory paid observation of an independent fair sign $E_0$, erased with probability $2\varepsilon$, $0\le\varepsilon\le1/2$, and a final paid audit of
\[
 (Z,E_0,s\delta_{\rm cal}),\qquad s\in\{-1,1\}.
\]
The fixed unknown sign $s$ affects no pre-audit law. The core, including its costs, has no access to $E_0$ beyond the initial visible tag. Use squared Euclidean loss on this product task, and charge these two additional calls. The core lower from \eqref{eq:central-lower}, evaluated with budget $N-2$ and safe projected dictionary width $M$, plus $2\varepsilon+\delta_{\rm cal}^2$, is a lower bound for either stopping interface.

Let $U(n,m)$ denote the upper expression in \eqref{eq:central-upper} for core budget $n$ and output width $m$, or its first-attempt alternative. For state stopping choose any $m\ge1$ with $1+3B+3m\le M$; for event stopping choose any $m\ge1$ with $\max(1+3B,3m)\le M$. Then an upper is
\begin{equation}\label{eq:joint-upper}
 U(N-2,m)+2\varepsilon+\delta_{\rm cal}^2.
\end{equation}
In the relative-order regular class, for $M$ large enough that these sufficient widths are a fixed fraction of $M$, this yields
\begin{equation}\label{eq:joint-law}
 R_{\mathrm{joint}}(N,M,\delta_{\rm cal},\varepsilon)
 \asymp F((N-2)/\kappa_-,M)+\delta_{\rm cal}^2+\varepsilon.
\end{equation}
The early sign channel has causal deficiency exactly $\varepsilon$ relative to the perfect-sign channel when the virtual sign must be supplied before any core call. The state counts are sufficient, not minimal for the enlarged experiment.
\end{theorem}
\begin{proof}
For the upper use one initial role, then three tagged copies of the core live roles for observed $+$, observed $-$ and erasure. Use three copies of its output dictionary. Decode the known sign or zero on erasure, and decode zero for the inaccessible last coordinate. This uses $1+3B$ live and $3m$ terminal roles, giving the stated widths by Lemma~\ref{lem:roles}. The common core upper is uniform after each tag. The other losses are exactly $2\varepsilon$ and $\delta_{\rm cal}^2$. The two compulsory calls give core budget $N-2$.

For the lower use the same actual core witness as in the central theorem. Conditioning on the initial tag does not change its full-report certificates; the initial sign and observer coins are additional independent randomization. The projected core decoder has at most $M$ values, so the actual-history proof gives its core lower. On erasure, induction over the core kernels and the stopping decisions leaves $E_0$ fair conditional on all available pre-audit information; its conditional squared prediction loss is at least one. Average the two indistinguishable fixed choices $s=\pm1$. For any last coordinate $r$, their mean loss is $\delta_{\rm cal}^2+r^2\ge\delta_{\rm cal}^2$. The core and erased-sign lower bounds hold under both of these identical pre-audit laws. Thus the three contributions hold simultaneously in one worst-case total task, rather than under incompatible separately optimized models. Regular geometry and fixed width ratios prove \eqref{eq:joint-law}.

On an erased branch every immediate virtual sign is independent of the true fair sign. The event that they agree has probability at most $1/2$, giving complete scored-law deficiency at least $\varepsilon$. Reporting the known sign otherwise and an independent fair sign on erasure gives a coupling with mismatch probability exactly $\varepsilon$. Run the same future core and audit in this coupling. The early-delivery requirement prevents waiting for the final audit to recover the sign, and proves the opposite inequality.
\end{proof}

\subsection*{A finite calibration acquisition with guarded stopping}
The reference and task-order certificates above are not inferred automatically from finite data. Here is one precise finite-sample certificate and its cost when the unknown quantities are the fixed probabilities of a physical command sampler. It also explains why high confidence alone cannot certify almost-sure stopping.

\begin{proposition}[Finite command calibration and a counted guard]\label{prop:pilot}
Suppose a physical sampler has a fixed unknown distribution $(\lambda_j)_{j=1}^J$ throughout calibration and deployment, its draws are independent, and deployment draws are independent of the current hidden input. The known target probabilities are $1/J$. Observe $n$ independent calibration draws at cost $c_{\mathrm{cal}}$ each, and let $\widehat\lambda_j$ be their frequencies. Fix $0<\zeta<1$, and choose a certified rational bound
\[
 r_n\ge\sqrt{\frac{\log(2J/\zeta)}{2n}}.
\]
Declare a relative certificate with parameter $\eta<1$ only when all intervals $[\widehat\lambda_j-r_n,\widehat\lambda_j+r_n]$ lie in $[(1-\eta)/J,(1+\eta)/J]$. The probability of a falsely declared certificate is at most $\zeta$.

On certificate failure, commit by a bounded-loss first-attempt fallback. On certificate success, deploy the corresponding upper of Theorem~\ref{thm:transfer}, with design hazard $1/b$, but force commitment with a fixed decoder after at most $T\ge1$ productive attempts. The whole-attempt cost bound $\kappa_+$ is required whether or not the certificate is true. If the true sampler lies in a tighter relative band $\eta_0<\eta$ and $2r_n\le(\eta-\eta_0)/J$, then the entire guarded procedure commits almost surely and has
\begin{align}
 \E C_{\mathrm{total}}&\le nc_{\mathrm{cal}}+\kappa_+(b+\zeta T),\label{eq:pilot-cost}\\
 R_{\mathrm{total}}&\le U+L\bigl(2\zeta+e^{-T/b}\bigr),\label{eq:pilot-risk}
\end{align}
where $U$ is the certified core upper. The calibration phase can be implemented with at most $(n+1)^J$ count-vector labels. Deployment uses a counted attempt index, multiplying its live-role bound by at most $T$; fallback and terminal roles are additional explicit roles. These resources may not be omitted from an autonomous budget.
\end{proposition}
\begin{proof}
For a Bernoulli variable $X$, the log moment-generating function of $X-\E X$ has second derivative at most $1/4$ under every exponential tilt, and value and derivative zero at the origin. Twice integrating gives $\log\E e^{t(X-\E X)}\le t^2/8$. Independence, the exponential Markov inequality and optimization in $t$ imply $\Prb(|\widehat\lambda_j-\lambda_j|>r)\le2e^{-2nr^2}$. A union bound proves that all confidence intervals contain the true probabilities except on an event of probability at most $\zeta$. Hence a false certificate has probability at most $\zeta$. If the tighter-band condition holds, this same simultaneous event puts each empirical interval within the declared wider band, so rejection of a true certificate also has probability at most $\zeta$.

On a true certificate, the conditional hazard proof bounds survival through $T$ attempts by $(1-1/b)^T\le e^{-T/b}$. The loss contributed by earlier acceptances is at most $U$; the forced last commitment adds at most $Le^{-T/b}$. Expected attempts on this branch are at most $b$. On a certificate failure the first-attempt fallback costs at most $\kappa_+$ and loses at most $L$. On a falsely accepted certificate the deterministic guard uses at most $T$ attempts and has loss at most $L$, even if the unguarded acceptance set has zero actual mass. Every attempt terminates almost surely by its finite whole-attempt cost bound. Combining these cases gives \eqref{eq:pilot-cost} and \eqref{eq:pilot-risk}, with the non-false cases bounded by $b\ge1$ attempts and the false-acceptance case by $\zeta T$ additional attempts.

A count vector stores every calibration total, including the number of draws already taken, and has at most $(n+1)^J$ possible values. At total $n$, the threshold comparison may be performed in charged scratch using a fixed certified finite arithmetic implementation; the resulting branch choice is encoded in the selected live role. The counts can then be discarded because the deployed dictionary is fixed for the nominal law and declared band, not fitted to the realized frequencies. The deployment guard requires its own retained index with $T$ values. This gives the asserted conservative finite-state accounts. Without the tighter-band assumption one still has the false-certificate probability bound and guarded termination, but no small rejection-risk bound; we do not omit that distinction.
\end{proof}

This result controls one genuinely acquired calibration quantity, not arbitrary unknown-kernel learning. It counts pilot observations, a physical fallback, the stopping guard and its memory. It prevents a rare incorrect calibration certificate from being hidden inside an unjustified robust almost-sure stopping assertion.

\appendix
\section{A predictive quotient constructed from positive-evidence histories}\label{sec:quotient}
The squared-task coordinates in \eqref{eq:derived} need not separate all future experiments. We therefore give the predictive realization independently of those coordinates, for a fixed prepared law with its declared legal test class. The theorem is a sufficient realization result and a detailed application of compact selection; predictive minimality itself has classical antecedents \cite{SC,BC}.

Let $\mathsf H$ be the standard Borel space of finite visible histories on which specified conditional versions are defined, embedded in a metrizable history carrier. Reports and resource records are included in the history. Actions belong to a finite set, with Borel legal domains. Histories with different legal continuations must be separated. Let $q_j(h)\in[0,1]$ be a countable determining list of conditional response coordinates for all finite open-loop future report words and bounded tests. Include legal-domain indicators and the coordinates necessary to determine report and cost-record laws. The list is closed in the following determining sense: the one-call prefixes of these future tests, localized by a countable measure-determining family of current-report functions, are determined by the list. This closure is a property of the declared executable tests, not a finite-moment assertion. Set $q(h)=(q_j(h))_{j\ge1}$ in $[0,1]^{\mathbb N}$.

\begin{theorem}[Compact-exhausted executable predictive realization]\label{thm:quotient}
Suppose $\mathsf H=\bigcup_{n\ge1}C_n$ with $C_n$ compact in its metrizable carrier and every $q_j$ continuous on each $C_n$. Assume the specified response coordinates are the consistent conditional laws of the raw causal kernels and satisfy the determining and continuation conditions above. Then the actual image $S=q(\mathsf H)$ is Borel and has a Borel section $s:S\to\mathsf H$. Equality of $q$ is exactly equality of all declared finite executable future laws, including legal domains. The report and next-quotient law factor through $q$ as a Borel kernel. This factorization is modulo the common current-report null sets of each compared pair of histories; it is not a pointwise claim on impossible reports.

If another Borel statistic $T$ determines the response coordinates by Borel maps $q_j=g_j\circ T$, then $q=g\circ T$ measurably, where $g=(g_j)_{j\ge1}$ on the realized range. Thus $q$ is minimal for these executable tests. Equality extends to bounded stopped tests which terminate almost surely under both compared histories, by finite truncation. Unbounded costs and infinite-horizon limits require their separate integrability or convergence hypotheses.
\end{theorem}
\begin{proof}
The countable product is compact metrizable, and coordinate continuity makes $q$ continuous on every $C_n$. Thus $q(C_n)$ is compact and $S=\bigcup_nq(C_n)$ is Borel. A section can be constructed rather than assumed. Given $x\in S$, first choose the least $n$ for which $x\in q(C_n)$, a Borel decision. For this compact $C_n$, choose finite closed covers at each scale, refining each previously chosen compact cell, with diameters tending to zero. Empty cells are discarded. At each level choose the first child compact cell $D$ for which $x\in q(D)$. Such a child exists because its parent covers the nonempty fiber; the test is Borel because $q(D)$ is compact. The chosen cells are nested compact sets with diameters tending to zero and each intersects the fiber. Choose a fixed representative in each possible cell. Their sequence is Borel in $x$ and converges to a unique point of $C_n$. Continuity of $q$ on $C_n$ shows that this point maps to $x$. The countable piecewise construction in the least $n$ gives the required Borel section.

Suppose $q(h)=q(h')$. The legal-domain coordinates give the same permitted next commands. The determining current-report coordinates give the same next-report law $R_a$ for a common legal command $a$. For a future response coordinate $q_j$ and a member $g$ of the determining current-report function family, continuation consistency gives a conditional joint-test value
\[
 \int g(y)q_j(hay)R_a(dy).
\]
The prefix closure and equality of $q(h),q(h')$ imply equality of this integral with the expression using $h'ay$. Uniqueness of finite measures on a determining family implies $q_j(hay)=q_j(h'ay)$ for $R_a$-almost every $y$. Intersecting these sets over the countable coordinates gives $q(hay)=q(h'ay)$ almost surely. Concatenations outside the positive-evidence history domain may be completed arbitrarily, because they have zero report mass from the histories in question.

Define the quotient report/update kernel by running the raw row from $s(x)$ and pushing the continued history through $q$. The section, the raw kernel and concatenation map are Borel, so this is a Borel kernel. The preceding almost-sure equality shows that it agrees with the pushforward from every representative of $x$, not just from $s(x)$. Induction gives finite-prefix equality under a common quotient policy. More generally finite adaptive policies depending on the visible future history are handled by the same raw conditional continuation laws: with finitely many actions, sum over action words and multiply their open-loop report laws by the measurable feedback probabilities of the policy. Equality extends from bounded determining functions to Borel sets by the monotone-class argument. The converse follows since each listed response is itself a declared test coordinate.

For another statistic with the stated Borel response maps, the product map $g$ is Borel and satisfies $q=g\circ T$; the asserted minimality is precisely this factorization, not a finite-state claim. For a stopped test truncate at time $n$ with a distinguished surviving outcome. Its finite law is determined. Almost-sure termination under both histories makes both surviving probabilities tend to zero, so the bounded stopped law agrees. This limit does not control expected stopping cost without an additional bound.
\end{proof}

\begin{proposition}[Primitive positive-evidence patches]\label{prop:raw-quotient}
The continuity and exhaustion conditions of Theorem~\ref{thm:quotient} hold for the following raw class. Hidden-state and report coordinates lie in compact metric spaces, the action set is finite, and primitive transition--report kernels have bounded continuous densities with respect to fixed finite reference measures; the preparation is fixed. A finite-dimensional positive instrument whose matrix entries depend continuously on compact report coordinates is another such class. Finite physical phases are included as discrete coordinates, with continuous legal-domain indicators on the corresponding patches. Only histories of positive prepared likelihood are included.
\end{proposition}
\begin{proof}
Fix a finite action word and history length. Iterated integration of products of bounded continuous primitive densities gives a continuous nonnegative history likelihood $p(h)$. The same statement holds for the numerator of a conditional finite future test whose report test function is continuous; the integration domains are fixed compact spaces with finite measures. Uniform boundedness on each finite product and dominated convergence justify continuity. In the matrix case these quantities are traces of products of continuous instrument maps followed by finite integration, with the same conclusion.

The sets where $p(h)\ge1/n$ are compact and exhaust the positive-likelihood histories for this action word. On each, division by $p(h)$ is continuous and all conditional response coordinates are continuous. Countably many lengths, finitely many action words per length, and countably many $n$ give the compact exhaustion. On a compact metric report space there is a countable uniformly dense family of continuous functions; include their finite products and all finite action words. Their integrals determine finite report laws, and adjoining finite prefixes gives the required determining closure. Costs and legal phases must be included among the stipulated continuous or discrete observable coordinates. The probability of a zero-likelihood report history is zero under the prepared law. The construction does not claim that the likelihood has a uniform positive lower bound over the whole experiment.
\end{proof}

The proposition is about continuity of raw finite-word responses. It does not assume a pre-existing posterior manifold or a covariance geometry, and it does not give a finite algorithm for discovering all quotient classes. An unknown-model family of separate posterior laws is not automatically one executable posterior. Applying this appendix uniformly to a model family would require a common determining response construction and compatible null-set/regularity conditions; the finite-state upper in the main theorem instead obtains its uniformity from the explicitly checked task-order program and stores no such family.

\section{Sharp boundaries of the proved comparisons}\label{sec:boundaries}
The following witnesses isolate distinct hypotheses. They are mathematical counterexamples to specific uncharged inferences, not a replacement of the main theorem by an impossibility statement.

\begin{proposition}[Absolute task error can be amplified by selection]\label{prop:absolute-sharp}
For every $b\ge1$ and $0\le e\le1/b$, there are positive prepared and actual joint laws with identical report marginals and one-call absolute squared-task defect at most $e$, whose optimal stopped risks at expected call budget $b$ differ by $be$, even with one output value. The difference occurs in either direction. A fixed relative task order is therefore strictly stronger than a small absolute task defect near a zero-loss component.
\end{proposition}
\begin{proof}
Draw an independent uniform report $x\in[0,1]$ at every call and let $D$ have mass $1/b$. All audits have conditional mean zero. Off $D$, both laws have a fresh fair-sign audit of magnitude one. On $D$, the comparison has a fair sign of magnitude $\sqrt{be}$ and the actual law has audit zero. For every $q\in[0,1]$ and finite-image readout $r(x)$, their unnormalized squared losses differ by $be\int_Dq(x)dx\le e$. The comparison residual variance is everywhere at least $be$, so every committing strategy has risk at least $be$; selecting $D$ and outputting zero attains it with expected cost $b$. The actual law attains zero by the same procedure. Reversing the laws proves the other direction. A single event-stopped label can reject at itself and commit with a final event; a state-stopped realization uses an additional live label. On $D$ a task zero under one law is positive under the other, so no two-sided relative order with a positive lower factor can hold.
\end{proof}

A report-only comparison also fails: a constant mark can have the same marginal under two laws while an additional visible bit reveals a fair audit sign. A two-output decoder then has zero loss, although a profile formed from the constant mark alone has residual variance one. The full-report conditional moments in \eqref{eq:full-moments} exclude this omission.


There is also an all-program witness for the necessity of more than report balance. Take the three admitted models $\theta\in\{-1,0,1\}$, let every precommit report be the same constant, and let the unseen audit at commitment be $Z=\theta$. There is one legal productive command, costing one, and no command changes $\theta$. The actual comparison member $\theta=0$ has $\Q_m(b)=0$. Nevertheless the distribution of every committed readout $R$ is identical under $\theta=\pm1$, because all precommit laws, including stopping and coins, are identical. The average of their risks is $1+\E R^2\ge1$, and output zero attains worst-case risk one. This holds for every almost-surely terminating policy, at every feasible budget and every dictionary size. Report masses match exactly; the missing first task moment is the obstruction. It is not repaired by choosing a different causal rejection rule.

A live-state count alone cannot justify a terminal-width subtraction. In the two-port interface, use labels $0,1$ with live commands gate and productive respectively, and decoder values $-1,+1$. Gate failure continues at $0$, success at $1$; a rejected productive report returns to $0$. A productive commitment emits a final event with label $0$ or $1$ according to its chosen output. Two live and two output roles have used only two labels. The state-stopped model cannot share those roles because its halt map is intrinsic. This is the exact semantic repair retained from the earlier work.

Finally, bounded-law discrepancy is not a cost certificate. A source procedure can agree with a target terminal output with probability one while performing an independent preliminary waiting time of arbitrarily large finite mean, or of infinite mean. Terminal total variation is then zero and expected cost is uncontrolled. Our simulator definition includes the scored stopping information when its law is compared, and still requires a separate expected-cost certificate because time is unbounded. Likewise, neither fixed-time equivalence nor finite-dimensional positive matrix identities imply large-deviation equivalence, unrestricted optimal-control equivalence, or closure of unbounded operators.

\subsection*{What is and is not generalized}
The central theorem generalizes exact and absolute-calibrated causal profile transfer to a common affine task order and derives its all-scale relative and stratified consequences. The positive-matrix characterization concerns localized task order; the quotient construction concerns compact-exhausted continuous-response experiments; the pilot concerns a fixed independent command sampler. These results do not establish necessity of a common acquisition program for every causal experiment, nor the minimal possible simulator state under all implementations, nor matching lower bounds for program length, scratch and arithmetic time. Those stronger mathematical questions are distinct from the fully stated sufficient-class theorem proved here. The general theorem does not depend on a conclusion taken from either realization. The retained nonintervening sensing, occupation-flow, exact role and absolute-defect results remain substantive independent parts of the complete companion.

\begin{thebibliography}{99}
\bibitem{Blackwell} D. Blackwell, \emph{Equivalent comparisons of experiments}, Ann. Math. Statist. \textbf{24} (1953), 265--272. DOI: 10.1214/aoms/1177729032.
\bibitem{CGM97} J. A. Cuesta-Albertos, A. Gordaliza and C. Matr\'an, \emph{Trimmed $k$-means: an attempt to robustify quantizers}, Ann. Statist. \textbf{25} (1997), 553--576. DOI: 10.1214/aos/1031833664.
\bibitem{GG99} L. A. Garc\'ia-Escudero and A. Gordaliza, \emph{Robustness properties of $k$ means and trimmed $k$ means}, J. Amer. Statist. Assoc. \textbf{94} (1999), 956--969. DOI: 10.1080/01621459.1999.10474200.
\bibitem{SC} C. R. Shalizi and J. P. Crutchfield, \emph{Computational mechanics: pattern and prediction, structure and simplicity}, J. Statist. Phys. \textbf{104} (2001), 817--879. DOI: 10.1023/A:1010388907793.
\bibitem{BC} N. Barnett and J. P. Crutchfield, \emph{Computational mechanics of input--output processes: structured transformations and the $\varepsilon$-transducer}, J. Statist. Phys. \textbf{161} (2015), 404--451; arXiv:1412.2690.
\bibitem{AMTW} A. Ambainis, M. Mosca, A. Tapp and R. de Wolf, \emph{Private quantum channels}, Proc. 41st Annual Symposium on Foundations of Computer Science (2000), 547--553. DOI: 10.1109/SFCS.2000.892142.
\bibitem{Buscemi} F. Buscemi, \emph{Comparison of quantum statistical models: equivalent conditions for sufficiency}, Comm. Math. Phys. \textbf{310} (2012), 625--647. DOI: 10.1007/s00220-012-1421-3.
\bibitem{Companion} Q. Qi, \emph{General Theta Foundations I: Acquired Geometry and Causal Resource Transfer}, complete technical companion on absolute calibration, exact stopping resources and nonintervening sensing, 2026. Source and independently built companion accompany this manuscript.
\end{thebibliography}

\end{document}
